% =============================================================================
% Scenario 2: Localized Crowd-Driven Service Degradation
% =============================================================================
% Include this section in your main LaTeX document.
% Prerequisites: booktabs, multirow, amsmath, hyperref packages.
% =============================================================================

\subsection{Scenario~2: Localized Crowd-Driven Service Degradation}
\label{sec:scenario_crowd_event}

\subsubsection{Motivation and Real-World Phenomenon}

Large-scale public gatherings---such as football matches, music concerts, political demonstrations, and religious festivals---constitute one of the most demanding operational scenarios for Mobile Edge Computing (MEC) infrastructures. During these events, tens of thousands of individuals converge on a single geographic location within a narrow time window, each carrying one or more network-connected devices that simultaneously consume and produce data-intensive content. Empirical measurements from deployed cellular networks have consistently demonstrated that such spatially concentrated crowds generate traffic volumes that exceed the capacity of local base stations by factors of 10$\times$ to 50$\times$ compared to normal operating conditions~\cite{Pujol2015Understanding, Shafiq2013First, Zhang2022Optimal}.

The phenomenon is characterized by several interrelated dynamics that distinguish it from ordinary urban traffic patterns. First, the \emph{spatial concentration} of demand is extreme: whereas everyday traffic is distributed across a metropolitan area, event-driven traffic is funnelled through a small number of cell towers and edge nodes located in the immediate vicinity of the venue. Second, the \emph{temporal dynamics} are highly bursty and correlated: significant moments during the event---a goal in a football match, the start of a headlining act, or a police confrontation at a protest---trigger synchronized behavioral responses in which a large fraction of the crowd simultaneously initiates video uploads, social media posts, or live-streaming sessions~\cite{Crane2008Robust}. Third, the \emph{application workload} is dominated by a small number of services: typically an official event application, a video streaming platform (e.g., YouTube, Twitch, or Instagram Live), and a social media or messaging service (e.g., Twitter/X, WhatsApp). Unlike the general-purpose traffic mix observed during normal urban operation, the crowd event workload is narrowly concentrated across two or three dominant applications.

From the perspective of service placement optimization, this scenario poses a particularly acute challenge: the surge in demand is \emph{localized} to a small subset of edge and fog nodes, it is \emph{temporally unpredictable} at fine granularity, and it generates \emph{cascading failures} as overloaded nodes degrade, edge links congest, and the placement solver must continuously re-optimize allocations under rapidly changing conditions. Understanding how placement algorithms respond to such extreme, localized pressure is essential for evaluating their robustness in real-world deployments~\cite{Armeniakos2025Stochastic, Fernando2025On}.

\subsubsection{Infrastructure Model}

The simulated infrastructure consists of $N = 50$ nodes organized in a composite Cloud--Fog--Edge hierarchy (\texttt{multi\_tier} topology), consistent with the reference architectures described in~\cite{Armeniakos2025Stochastic, Fernando2025On}. The hierarchical structure is generated as follows:

\begin{itemize}
    \item \textbf{Cloud layer:} Approximately 2\% of nodes ($\sim$1 node) are designated as cloud data centers. These nodes possess virtually unlimited computational resources (100\,000~units of RAM and CPU), reflecting the elastic capacity of centralized cloud infrastructure, but are located at higher network distances from the edge.
    
    \item \textbf{Fog layer:} 10\% of nodes ($\sim$5 nodes) serve as fog gateways, each connected to 2 edge nodes upon creation ($m = 2$). These intermediate nodes provide higher capacity than edge devices but are geographically closer to the venue.
    
    \item \textbf{Edge layer:} The remaining $\sim$44 nodes constitute the edge tier, connected via a random geometric graph with radius $r = 0.15$. This radius is slightly larger than in the baseline scenario to ensure full network connectivity despite the tight spatial clustering of users.
    
    \item \textbf{Layer assignment:} Nodes are classified into layers based on betweenness centrality using a centrality-based assignment policy. Nodes with centrality above 0.10 are promoted to the cloud layer, those above 0.02 to the fog layer, and the remainder constitute the edge tier.
\end{itemize}

\paragraph{Node resources.} RAM and CPU capacities are assigned using a Pareto distribution with shape parameter $a = 1.16$, directly proportional to each node's centrality within the network graph. Edge nodes range from 4~GB to 32~GB of RAM and 2 to 16 CPU cores, following the heavy-tailed resource distribution observed in production data center traces~\cite{Reiss2012Heterogeneity}. This resource heterogeneity is critical for this scenario: the few edge nodes closest to the venue must serve a disproportionate share of the crowd's demand, and their limited capacity makes them vulnerable to saturation.

\paragraph{Node events.} Hardware failures are modeled as rare events with an exponential inter-arrival time of $\lambda = 1/500$ (mean interval of 500 time units between failures), consistent with the failure rates reported in~\cite{Gill2011Understanding}. Recovery times follow a log-normal distribution ($\mu = 2.0$, $\sigma = 0.8$). In contrast, \emph{node degradation} events are extremely frequent ($\lambda = 1/30$), modeling the thermal throttling, memory pressure, and operating system contention that occur when edge servers are overwhelmed by concentrated demand~\cite{Bantock2020Mitigating, Xu2018Fail}. Each degradation event reduces both RAM and CPU capacity by a fraction drawn from a Beta distribution with parameters $\alpha = 3$, $\beta = 4$ (mean loss $\approx 43\%$), substantially more severe than in baseline conditions because all user traffic is funnelled through a small number of nearby nodes. Recovery from degradation follows an exponential distribution with scale~25, reflecting the sustained nature of the crowd-induced pressure.

\paragraph{Edge attributes and events.} Link delays are drawn from a log-normal distribution ($\mu = 1.5$, $\sigma = 0.8$), consistent with the measurements reported by Hao et al.~\cite{Hao2010Optimal}. Link failures are rare ($\lambda = 1/400$), but \emph{edge congestion} events occur with very high frequency ($\lambda = 1/20$). When a congestion event fires, the affected link suffers bandwidth loss drawn from a Beta($\alpha = 4$, $\beta = 3$) distribution (mean loss $\approx 57\%$) and delay inflation drawn from a log-normal distribution ($\mu = 0.8$, $\sigma = 0.3$, median multiplier $\approx 2.2\times$). These values reflect empirical observations of cellular link saturation during mass events~\cite{Shafiq2013First, Zhang2022Optimal}. Congestion clearance follows an exponential distribution with scale~15, modeling the sustained nature of crowd-driven traffic.

\subsubsection{Application Workload}

The simulation is initialized with exactly $K = 3$ applications, reflecting the empirical observation that crowd events are dominated by a small number of services~\cite{Pujol2015Understanding}. These three applications persist throughout the entire simulation duration---new application creation and application removal are effectively disabled (exponential inter-arrival time with scale~$99\,999$, yielding a negligible probability of occurrence during the 5\,000-iteration simulation horizon). This design choice models the realistic scenario in which attendees at a stadium or concert use only the official event application, a video streaming platform, and a social media or messaging service.

\paragraph{Service profiles.} Each application's internal topology follows a directed scale-free graph model, with 2 initial nodes and 2 edges per new node, generating the hub-dominated dependency structures typical of real microservice architectures. The number of constituent services per application ranges from 3 to 8. Three service profiles are defined:

\begin{itemize}
    \item \textbf{Video streaming (55\%):} Representing live broadcasts, video calls, and short-form video uploads (TikTok, Instagram Stories). CPU requirements follow a log-normal distribution ($\mu = 2.5$, $\sigma = 1.0$) and RAM requirements ($\mu = 3.5$, $\sigma = 1.5$), reflecting the compute-intensive nature of real-time video encoding and transcoding~\cite{CiscoAIR2023}.
    
    \item \textbf{Social media (30\%):} Representing photo sharing, live tweeting, and group messaging. Resource requirements are moderate: CPU ($\mu = 1.0$, $\sigma = 0.8$) and RAM ($\mu = 1.5$, $\sigma = 1.0$). The proportion of social media services is higher than in baseline scenarios because crowd events drive massive social media activity~\cite{Benevenuto2009Characterizing}.
    
    \item \textbf{IoT/lightweight (15\%):} Representing auxiliary services such as venue sensors, ticketing systems, or wearable device telemetry. These require minimal resources: CPU ($\mu = 0.2$, $\sigma = 0.5$) and RAM ($\mu = 0.2$, $\sigma = 0.5$).
\end{itemize}

\paragraph{Network requirements.} Applications impose latency constraints drawn from $\mathcal{U}(5, 30)$~ms, reflecting the tight latency budgets required for real-time streaming and augmented reality overlays~\cite{3GPP2018System}. Bandwidth requirements range from 20 to 150~Mbps, consistent with the traffic volumes reported in the Cisco Annual Internet Report for video-dominated mobile scenarios~\cite{CiscoAIR2023}.

\paragraph{Application popularity.} Initial user-to-application assignment follows a Zipf distribution with exponent $a = 1.2$. A moderate exponent is chosen because, with only three applications, a highly skewed distribution (e.g., $a \geq 2.0$) would render two of the three applications nearly unused, which is inconsistent with the empirical observation that all major services see substantial usage during crowd events~\cite{Breslau1999Web, Pujol2015Understanding}. Each application is initially assigned between 8 and 12 users, yielding a total initial population of approximately 30 users.

\paragraph{Application dynamics.} Despite the absence of application creation and removal, the three resident applications undergo frequent dynamic changes:

\begin{itemize}
    \item \emph{Footprint updates} ($\lambda = 1/50$): Auto-scaling events that adjust CPU and RAM requirements by a factor drawn from $\mathcal{U}(0.10, 0.25)$, modeling the elastic resource provisioning that occurs under extreme localized load~\cite{Cortez2017Resource}.
    
    \item \emph{Network requirement updates} ($\lambda = 1/60$): Bandwidth and latency constraint adjustments ($\mathcal{U}(0.08, 0.20)$) reflecting the fluctuating network demands driven by user count variations.
    
    \item \emph{Popularity surges} ($\lambda = 1/35$): Very frequent events modeling viral moments---goals, key speeches, or confrontations---that cause a sudden spike in one application's popularity. 85\% of surges are transient with duration drawn from a Weibull distribution ($a = 1.8$, $\lambda = 8.0$)~\cite{Crane2008Robust}. The remaining 15\% are permanent, modeling sustained shifts in usage patterns (e.g., a second-half formation change that drives sustained tactical discussion).
    
    \item \emph{Popularity drops} ($\lambda = 1/80$): Less frequent than surges, reflecting the observation that crowd interest is sustained during the event. 90\% of drops are transient (Weibull, $a = 1.5$, $\lambda = 8.0$).
    
    \item \emph{Geographic demand shifts} ($\lambda = 1/15$): Very frequent hotspot position adjustments modeling crowd movement within the venue (halftime, entrance/exit waves, route changes for mobile protests). Demand centers shift via a combination of jump relocation (probability~0.3) and Gaussian drift ($\sigma = 0.10$)~\cite{Camp2002Survey, Wang2019Delay}.
\end{itemize}

\subsubsection{User Population and Mobility}

\paragraph{Spatial distribution.} Users are generated according to a Thomas cluster process with a single hotspot ($H = 1$) and an extremely tight spatial spread ($\sigma_s = 0.015$). This configuration concentrates the entire user population within a narrow radius around the venue---the stadium, concert hall, or protest route---reflecting the defining characteristic of crowd events. The spread parameter is tighter than in multi-hotspot scenarios (which typically use $\sigma_s \geq 0.02$) because crowd events, by definition, force all participants into a single geographic area~\cite{Saha2016Enriched}.

\paragraph{Initial population.} The simulation begins with $U = 30$ users, representing 60\% of the number of infrastructure nodes. This ratio is intentionally higher than in baseline scenarios (which typically use 30--40\% of node count) to model the large initial crowd already present when the event begins.

\paragraph{Mobility model.} User movement follows a Manhattan mobility model on a $15 \times 15$ grid spanning a $2.0 \times 2.0$ spatial region. Turn probabilities are set to $p_{\text{straight}} = 0.40$, $p_{\text{left}} = 0.30$, $p_{\text{right}} = 0.30$, reflecting the frequent direction changes observed in crowd settings where attendees mill around the venue, visit concession stands, and navigate corridors. Movement speed is drawn from $\mathcal{U}(0.015, 0.040)$, approximately twice the speed used in normal urban scenarios ($\mathcal{U}(0.008, 0.025)$), consistent with pedestrian speed measurements in moderate crowd densities ($\sim$1.5--2.5~m/s)~\cite{Weidmann1993Transporttechnik, Helbing2000Simulating}. Movement events fire with high frequency ($\lambda = 1/20$), creating a highly dynamic spatial distribution in which users continuously change their associated edge node.

\paragraph{User request behavior.} Each user generates service requests at a rate drawn from an exponential distribution with scale~12, higher than the baseline value of 8, reflecting the elevated per-user activity observed at live events where all attendees are actively using their devices for streaming, messaging, and social media~\cite{Pujol2015Understanding}. Request ratio changes occur with frequency $\lambda = 1/40$ and are biased toward increases (70\% probability), with multipliers drawn from $\mathcal{U}(2.5, 5.0)$, modeling the demand spikes triggered by exciting moments during the event. Restoration of request rates follows a log-normal distribution ($\mu = 0.8$, $\sigma = 0.4$).

\paragraph{Fixed user lifetime.} A distinguishing feature of this scenario is the deterministic user lifetime. Each user's removal event is scheduled according to a degenerate uniform distribution $\mathcal{U}(300, 300)$, guaranteeing that every user persists for exactly 300 time units before being removed from the simulation. This design choice models the fixed duration of the public event: all attendees arrive, participate for the duration of the event, and depart when it concludes. Upon removal, all associated pending events (movement, request ratio changes, suspensions) are automatically cleaned up by the simulator's event garbage collection mechanism.

\paragraph{User disconnections.} Temporary user suspensions (dropped connections) occur with frequency $\lambda = 1/80$, slightly more frequent than in baseline scenarios due to the extreme network congestion near the single hotspot. Reconnection times follow a log-normal distribution ($\mu = 0.8$, $\sigma = 0.4$)~\cite{Haibeh2022A}.

\subsubsection{Crowd Surge: Composed Event Mechanism}

The central dynamic element of this scenario is the \emph{crowd surge} composed event, which models the synchronized behavioral response of a large crowd to a significant moment during the event. Crowd surges fire with an exponential inter-arrival time of scale~80, corresponding to approximately 10--15 surges during a typical 5\,000-iteration simulation. This frequency is consistent with the number of ``exciting moments'' observed during real-world events: a 90-minute football match typically contains 10--15 significant play events that trigger coordinated audience reactions~\cite{Pujol2015Understanding}.

Each crowd surge simultaneously triggers four correlated sub-actions:

\begin{enumerate}
    \item \textbf{Mass user creation (20 new users):} A burst of 20 new users is instantaneously injected into the simulation, modeling a wave of attendees who activate their devices or new arrivals entering the venue. Each new user inherits the same spatial distribution (single hotspot, $\sigma_s = 0.015$), mobility parameters (Manhattan model, $\mathcal{U}(0.015, 0.040)$ speed), and deterministic lifetime (300 time units) as the initial population. Target selection follows a random strategy over the existing user population, but the \texttt{new\_user} action ignores the target identifier and creates a fresh user entity, ensuring that each sub-action produces an independent new user.
    
    \item \textbf{Application popularity surge (3 applications):} The popularity ranking of 3 applications is simultaneously surged via the \texttt{surge\_popularity} action. Target selection uses a \texttt{random\_proximity} strategy that preferentially selects applications whose users are located near the surge epicenter. Surges are 90\% transient, with duration drawn from a Weibull distribution ($a = 2.0$, $\lambda = 15.0$), and 10\% permanent.
    
    \item \textbf{Edge congestion (4 links):} Four network links are simultaneously congested, modeling the traffic spike at cell towers and backhaul links serving the venue area. Bandwidth loss follows a Beta($\alpha = 4$, $\beta = 2$) distribution (mean loss $\approx 67\%$), substantially more severe than steady-state congestion, because crowd surge moments trigger synchronized uploads and streams from a large fraction of the audience. Delay inflation follows a log-normal distribution ($\mu = 1.0$, $\sigma = 0.3$, median $\approx 2.7\times$). Congestion clearance follows an exponential distribution with scale~20.
    
    \item \textbf{Request ratio spike (8 users):} Eight users experience an immediate increase in their service request rate, with multipliers drawn from $\mathcal{U}(3.0, 6.0)$. The increase probability is set to 1.0 (always increase during a surge), and restoration follows a log-normal distribution ($\mu = 0.5$, $\sigma = 0.3$), reflecting the rapid decay of excitement-driven activity.
\end{enumerate}

The combination of these four sub-actions creates a realistic, multi-dimensional shock to the MEC infrastructure: new users appear (increasing placement demand), applications become more popular (shifting the optimization landscape), network links degrade (constraining the feasible placement space), and individual user demand spikes (increasing per-user resource consumption). This composed event mechanism allows the simulation to capture the \emph{correlated, multi-modal} nature of crowd-driven disruptions, which is fundamentally different from the independent, single-attribute perturbations assumed in most existing benchmarks.

\subsubsection{Continuous User Arrivals}

In addition to the crowd surge bursts, the simulation includes a continuous stream of individual user arrivals ($\lambda = 1/15$), modeling the steady trickle of late-arriving attendees entering the venue. Each new user is created with the same configuration parameters described above. This steady-state arrival process ensures that the user population grows gradually between surges, providing a realistic baseline growth trajectory upon which the discrete crowd surge bursts are superimposed.

\subsubsection{Summary of Key Configuration Parameters}

Table~\ref{tab:scenario2_params} summarizes the principal configuration parameters for this scenario, along with their values, distributional forms, and the empirical references that motivated their selection.

\begin{table}[htbp]
\centering
\caption{Configuration parameters for Scenario~2 (Localized Crowd-Driven Service Degradation).}
\label{tab:scenario2_params}
\small
\begin{tabular}{@{}llll@{}}
\toprule
\textbf{Parameter} & \textbf{Value / Distribution} & \textbf{Rationale} \\
\midrule
\multicolumn{3}{@{}l}{\textit{Infrastructure}} \\
\quad Nodes & 50 (multi-tier) & \cite{Armeniakos2025Stochastic} \\
\quad Cloud / Fog / Edge & 2\% / 10\% / 88\% & \cite{Fernando2025On} \\
\quad RAM (edge) & Pareto($a{=}1.16$), 4--32\,GB & \cite{Reiss2012Heterogeneity} \\
\quad Node degradation freq. & Exp($\lambda{=}1/30$) & \cite{Bantock2020Mitigating} \\
\quad Node capacity loss & Beta($3, 4$), mean $\approx$43\% & \cite{Xu2018Fail} \\
\quad Edge congestion freq. & Exp($\lambda{=}1/20$) & \cite{Zhang2022Optimal} \\
\quad BW loss (congestion) & Beta($4, 3$), mean $\approx$57\% & \cite{Shafiq2013First} \\
\midrule
\multicolumn{3}{@{}l}{\textit{Applications}} \\
\quad Number of apps & 3 (persistent) & \cite{Pujol2015Understanding} \\
\quad Profiles & Video 55\%, Social 30\%, IoT 15\% & \cite{CiscoAIR2023} \\
\quad Popularity & Zipf($a{=}1.2$) & \cite{Breslau1999Web} \\
\quad Surge freq. & Exp($\lambda{=}1/35$) & \cite{Crane2008Robust} \\
\quad Geo-demand shift freq. & Exp($\lambda{=}1/15$) & \cite{Camp2002Survey} \\
\midrule
\multicolumn{3}{@{}l}{\textit{Users}} \\
\quad Initial users & 30 & 60\% of node count \\
\quad Hotspots & 1 ($\sigma_s{=}0.015$) & \cite{Saha2016Enriched} \\
\quad Speed & $\mathcal{U}(0.015, 0.040)$ & \cite{Weidmann1993Transporttechnik} \\
\quad Movement freq. & Exp($\lambda{=}1/20$) & \cite{Bettstetter2004Stochastic} \\
\quad Lifetime & Deterministic (300 t.u.) & Fixed event duration \\
\quad Request rate & Exp($\text{scale}{=}12$) & \cite{Pujol2015Understanding} \\
\midrule
\multicolumn{3}{@{}l}{\textit{Crowd Surge (composed event)}} \\
\quad Frequency & Exp($\lambda{=}1/80$) & \cite{Pujol2015Understanding} \\
\quad New users per surge & 20 & Burst of attendees \\
\quad Apps surged & 3 & Proximity-based \\
\quad Edges congested & 4 & Venue-local towers \\
\quad Users req. spiked & 8 & \cite{Crane2008Robust} \\
\quad BW loss (surge) & Beta($4, 2$), mean $\approx$67\% & Extreme simultaneous load \\
\quad Req. multiplier (surge) & $\mathcal{U}(3.0, 6.0)$ & Synchronized demand \\
\bottomrule
\end{tabular}
\end{table}
